\documentclass[11pt]{amsart}
\usepackage[T1]{fontenc}
\usepackage[utf8]{inputenc}
\usepackage{lmodern}
\usepackage{amsmath,amssymb,amsthm,mathtools}
\usepackage{microtype}
\usepackage{enumitem}
\usepackage[hidelinks]{hyperref}
\usepackage[a4paper,margin=32mm]{geometry}
\newtheorem{theorem}{Theorem}[section]
\newtheorem{lemma}[theorem]{Lemma}
\newtheorem{proposition}[theorem]{Proposition}
\newtheorem{corollary}[theorem]{Corollary}
\theoremstyle{definition}
\newtheorem{definition}[theorem]{Definition}
\newcommand{\Z}{\mathbb Z}
\newcommand{\Q}{\mathbb Q}

\newcommand{\R}{\mathbb R}
\newcommand{\F}{\mathbb F}
\newcommand{\Aut}{\operatorname{Aut}}
\newcommand{\Inn}{\operatorname{Inn}}
\newcommand{\Out}{\operatorname{Out}}
\newcommand{\Tor}{\operatorname{Tor}}
\newcommand{\rank}{\operatorname{rank}}
\newcommand{\core}{\operatorname{core}}

\title{Balanced two-generator $p$-groups of deficiency zero and unbounded nilpotency class}
\author{GroupTheory50}
\date{August 27, 2026}
\subjclass[2020]{20D15, 20F05, 20F18}
\keywords{finite p-groups, deficiency, nilpotency class, balanced presentation, Kourovka Notebook}

\begin{document}
\begin{abstract}
For every prime $p>3$ we construct an explicit family of finite two-generator $p$-groups of deficiency zero whose nilpotency classes tend to infinity.  The construction gives, for every $n\geq1$, a group of order $p^{2n+1}$ and class $n+1$.  This gives an affirmative answer to Problem 17.113 of the Kourovka Notebook.
\end{abstract}
\maketitle

\section{Introduction}
The deficiency of a finite presentation
\[
  \langle x_1,\dots,x_d\mid r_1,\dots,r_s\rangle
\]
is $d-s$, and the deficiency of a finitely presented group is the supremum of the deficiencies of its finite presentations.  A presentation with the same number of generators and relators is usually called balanced.  Balanced presentations of finite $p$-groups are constrained enough to be useful in topology and combinatorial group theory, but flexible enough to support groups of substantial nilpotency class.

Problem 17.113 of the Kourovka Notebook, proposed by J. Wiegold in the 17th issue, asks whether, for a fixed prime $p>3$, there are two-generator finite $p$-groups of deficiency zero with arbitrarily high nilpotency class \cite{kourovka}.  We give a uniform metacyclic construction.

\section{The family}
Fix a prime $p>3$ and an integer $n\geq1$.  Define
\[
 G_n=\left\langle a,b\ \middle|\ b^{-1}ab=a^{1+p},\quad b^{p^n}=a^{p^n}\right\rangle.
\]
We use the commutator convention $[x,y]=x^{-1}y^{-1}xy$.

\begin{theorem}\label{thm:main}
For every $n\geq1$, the group $G_n$ is a finite two-generator $p$-group of order $p^{2n+1}$, nilpotency class $n+1$, and deficiency zero.
\end{theorem}

\begin{proof}
Conjugating the relation $b^{p^n}=a^{p^n}$ by $b$ gives
\[
 b^{p^n}=a^{p^n(1+p)}.
\]
Comparison with the defining relation yields $a^{p^{n+1}}=1$, and then also $b^{p^{n+1}}=1$.  Using $b^{-1}ab=a^{1+p}$, every word collects to $a^ib^j$.  The relation $b^{p^n}=a^{p^n}$ allows us to choose
\[
 0\leq i<p^{n+1},\qquad 0\leq j<p^n,
\]
so $|G_n|\leq p^{2n+1}$.

For the reverse inequality, put $q=1+p$ and let
\[
 A_n=\langle\alpha\rangle\cong C_{p^{n+1}},\qquad
 B_n=\langle\beta\rangle\cong C_{p^{n+1}}.
\]
Let $B_n$ act on $A_n$ by $\alpha^\beta=\alpha^q$.  The standard lifting-the-exponent calculation gives
\[
 v_p(q^{p^k}-1)=k+1,
\]
so this automorphism has order $p^n$.  In $S_n=A_n\rtimes B_n$ the element
\[
 z=\beta^{p^n}\alpha^{-p^n}
\]
is central of order $p$.  Moreover $\langle z\rangle$ meets both $A_n$ and $B_n$ trivially.  Hence
\[
 H_n=S_n/\langle z\rangle
\]
has order $p^{2n+1}$.  The images of $\alpha$ and $\beta$ satisfy the defining relations of $G_n$, so there is an epimorphism $G_n\twoheadrightarrow H_n$.  The upper bound already obtained forces $G_n\cong H_n$.

The first defining relation gives
\[
 [a,b]=a^p,
\]
and therefore, for $k\geq0$,
\[
 [a^{p^k},b]=a^{p^{k+1}}.
\]
It follows inductively that
\[
 \gamma_r(G_n)=\langle a^{p^{r-1}}\rangle\qquad(2\leq r\leq n+2).
\]
Thus $\gamma_{n+1}(G_n)\neq1$ and $\gamma_{n+2}(G_n)=1$, so the class is $n+1$.

The displayed presentation has two generators and two relators, hence $\operatorname{def}(G_n)\geq0$.  On the other hand, if a finite group has a presentation with $d$ generators and $s$ relators, abelianization and tensoring the relation matrix with $\Q$ show that $s\geq d$; thus every finite group has deficiency at most zero.  Consequently $\operatorname{def}(G_n)=0$.
\end{proof}

\begin{corollary}\label{cor:unbounded-class}
By Theorem~\ref{thm:main}, for every prime $p>3$ there exist two-generator finite $p$-groups of deficiency zero and arbitrarily large nilpotency class.
\end{corollary}

\section{Further remarks}
The groups $G_n$ are metacyclic central quotients of semidirect products of two cyclic $p$-groups.  In particular, the growth of nilpotency class does not require a growing number of generators or relators.  The same presentation also gives a concrete family suitable for computational experiments with balanced finite presentations.

The restriction $p>3$ is the one appearing in the Kourovka problem.  The construction itself uses only the familiar $p$-adic behavior of $1+p$ for odd $p$; we have not attempted here to optimize the small-prime range.

\begin{thebibliography}{99}
\bibitem{kourovka} E. I. Khukhro and V. D. Mazurov (eds.), \emph{Unsolved Problems in Group Theory: The Kourovka Notebook}, No. 21, Novosibirsk, 2026; arXiv:1401.0300.
\bibitem{hall} M. Hall, Jr., The theory of groups, Macmillan, New York, 1959.
\bibitem{robinson} D. J. S. Robinson, \emph{A Course in the Theory of Groups}, 2nd ed., Springer, 1996.
\end{thebibliography}
\end{document}
